This section describes the image processing techniques used for tumor edge detection, following the methodology of Zotin et al. \cite{zotin2018edge}. The pipeline has three phases: pre-processing for noise reduction and contrast enhancement, clustering-based segmentation, and edge extraction. \gls{mri} scans are affected by noise and artifacts coming from the physical limits of the acquisition equipment. The methodology uses a Median Filter, a non-linear rank-order filter, to suppress impulse noise while preserving edge sharpness. Given a kernel $K$ centered at $(x,y)$, the operation is defined as

\begin{equation}
    I_{new}(x,y) = \text{med}_{(k_x,k_y) \in K} \{ I_{old}(k_x,k_y) \}
\end{equation}
To further highlight \glspl{roi}, the \gls{bcet} is applied. This technique expands the image contrast without altering the histogram pattern \cite{zotin2018edge} using a parabolic function
\begin{equation}
    I_{new} = a(I_{old} - b)^2 + c
\end{equation}
where coefficients are computed deterministically based on the maximum and minimum intensity values of the output image ($I_{new}$), and the mean value of intensity of the output image ($E$)
\begin{equation}
b = \frac{h^2(E - L) - s(H - L) + l^2(H - E)}
{2[h(E - L) - e(H - L) + l(H - E)]}
\end{equation}
\begin{equation}
a = \frac{H - L}{(h - l)(h + l - 2b)}
\end{equation}
\begin{equation}
c = L - a(l - b)^2
\end{equation}
\begin{equation}
s = 1/N \sum_{i=1}^{N}I_{old}^2(i)
\end{equation}
Segmentation aims to partition image pixels into distinct segments based on similarity. Unlike "hard" clustering, \gls{fcm} \cite{bezdek1981pattern} is a "soft" method that assigns a membership degree $u_{ik} \in [0,1]$ to each pixel. The algorithm iteratively minimizes the objective function
\begin{equation}
    F(W,V) = \sum_{i=1}^{c} \sum_{k=1}^{n} (u_{ik})^m \cdot ||x_k - v_i||^2
\end{equation}
where $ 1\le m \le \infty$ is the "fuzziness" index, $w_{ik}\in [0,1]$ is the degree of membership of $ x_k$ in the $i$th cluster, and $\|x_k - v_i\|^2$ represents the distance between the data $x_k$ and the cluster center $v_i$. This soft assignment is well suited to the Partial Volume Effect inherent to \gls{mri}: a single voxel at the boundary between two tissue types can physically contain a mixture of both, so a hard partition such as $k$-means is forced to assign it entirely to one class, while \gls{fcm}'s fractional membership reflects this underlying ambiguity directly, generally producing more anatomically coherent clusters at tissue boundaries. A simpler alternative, Otsu's global thresholding, is computationally cheaper and depends on a single, easily interpreted parameter, but it requires a histogram with well-separated bimodal peaks; \gls{mri} histograms are frequently smoothed by noise and tissue heterogeneity, in which case Otsu's single rigid boundary misclassifies voxels that \gls{fcm}'s gradual membership transition handles more robustly. Finally, an edge is a strong local change in brightness that usually corresponds to an anatomical boundary. The Canny algorithm is used here because it satisfies three criteria identified by Pratt \cite{pratt1979quantitative}: low error rate, precise localization, and a single-pixel response. The process has four stages:
\begin{enumerate}
    \item \textbf{Smoothing}: Convolution with a Gaussian filter to eliminate residual noise.
    \item \textbf{Gradient Calculation}: Determination of spatial gradient magnitude and direction.
    \item \textbf{Non-Maximal Suppression}: Thinning the gradient ridges to isolate the pixels with the highest intensity change.
    \item \textbf{Hysteresis Thresholding}: Validating edges using high and low thresholds to differentiate true anatomical boundaries from noise-induced segments.
\end{enumerate}